% =====================================================================
%  CARNET D'ENTRAÎNEMENT — FICHE 8 — THÈME 5
%  Tuto : Équilibre chimique et loi de Le Chatelier
% =====================================================================
\documentclass[11pt,a4paper]{article}
\usepackage[utf8]{inputenc}
\usepackage[T1]{fontenc}
\usepackage{lmodern}
\usepackage[french]{babel}
\usepackage{amsmath,amssymb,mathtools}
\usepackage{icomma}
\usepackage{siunitx}
\sisetup{output-decimal-marker={,},per-mode=power,group-separator={\,},group-minimum-digits=5}
\usepackage[version=4]{mhchem}
\usepackage{xcolor}
\usepackage{colortbl}
\usepackage{tikz}
\usepackage[most]{tcolorbox}
\usepackage{enumitem}
\usepackage{array}
\usepackage{booktabs}
\usepackage{fancyhdr}
\usepackage{caption}
\captionsetup{labelformat=empty,justification=centering,font=small}
\usepackage[hidelinks]{hyperref}
\usetikzlibrary{arrows.meta,positioning,calc}
\usepackage[a4paper,margin=2.0cm,headheight=26pt,headsep=10pt]{geometry}
\usepackage{titlesec}

\definecolor{primary}{RGB}{150,98,162}
\definecolor{primarydark}{RGB}{92,52,110}
\definecolor{defcol}{RGB}{0,114,178}
\definecolor{thmcol}{RGB}{0,140,120}
\definecolor{formulacol}{RGB}{198,93,9}
\definecolor{methcol}{RGB}{90,90,100}
\definecolor{attncol}{RGB}{200,40,55}
\definecolor{excol}{RGB}{0,135,90}
\definecolor{reglecol}{RGB}{150,98,162}

\tcbset{boxbase/.style={enhanced,breakable,boxrule=0.9pt,arc=2.5pt,left=3mm,right=3mm,top=2.3mm,bottom=2.3mm,fonttitle=\bfseries,coltitle=white,toptitle=0.6mm,bottomtitle=0.6mm}}
\newtcolorbox{definitionb}[1][]{boxbase,colback=defcol!5,colframe=defcol,title={\textsc{Définition}\ifx\relax#1\relax\relax\else\textmd{ --- \itshape #1}\fi}}
\newtcolorbox{formuleb}[1][]{boxbase,colback=formulacol!6,colframe=formulacol,title={\textsc{Formule}\ifx\relax#1\relax\relax\else\textmd{ --- \itshape #1}\fi}}
\newtcolorbox{theoremeb}[1][]{boxbase,colback=thmcol!6,colframe=thmcol,title={\textsc{Propriété}\ifx\relax#1\relax\relax\else\textmd{ --- \itshape #1}\fi}}
\newtcolorbox{attentionb}[1][]{boxbase,colback=attncol!6,colframe=attncol,title={\textsc{Attention}\ifx\relax#1\relax\relax\else\textmd{ : \itshape #1}\fi}}

\newcommand{\dHr}{\Delta H^{\circ}_{\text{r}}}

\titleformat{\section}{\large\bfseries\color{primarydark}}{\thesection}{0.6em}{}

\pagestyle{fancy}
\fancyhf{}
\fancyhead[L]{\small\textcolor{primarydark}{\textbf{Fiche 8} — Thème 5 : Équilibre chimique et Le Chatelier}}
\fancyhead[R]{\small\textcolor{primarydark}{\textsc{Tuto}}}
\fancyfoot[C]{\small\thepage}
\renewcommand{\headrulewidth}{0.8pt}
\renewcommand{\headrule}{\hbox to\headwidth{\color{primary}\leaders\hrule height \headrulewidth\hfill}}
\setlength{\parindent}{0pt}\setlength{\parskip}{2.5pt}

\begin{document}

\begin{center}
\begin{tikzpicture}
\node[rectangle,rounded corners=7pt,fill=primary,text=white,minimum width=16.4cm,
      minimum height=1.5cm,align=center,inner sep=8pt]
  {{\Large\bfseries Thème 5 — Équilibre chimique et loi de Le Chatelier}\\[2pt]
   {\small La constante $K$ et le déplacement d'équilibre}};
\end{tikzpicture}
\end{center}

\section{Équilibre et constante d'équilibre}

\begin{definitionb}[état d'équilibre]
Une réaction \textbf{réversible} ($\rightleftharpoons$) atteint l'\textbf{équilibre} quand la vitesse du
sens direct égale celle du sens inverse : les concentrations ne varient plus (état stationnaire), mais les
deux réactions continuent (équilibre \emph{dynamique}). « Équilibre » ne veut \emph{pas} dire « quantités égales ».
\end{definitionb}

\begin{formuleb}[constante d'équilibre (loi d'action de masse)]
Pour $\;a\,\ce{A} + b\,\ce{B} \rightleftharpoons c\,\ce{C} + d\,\ce{D}\;$ (concentrations à l'équilibre) :
\[ \boxed{\;K = \dfrac{[\ce{C}]^{c}\,[\ce{D}]^{d}}{[\ce{A}]^{a}\,[\ce{B}]^{b}}\;} \qquad
   \text{\small produits au numérateur, exposants = coefficients.} \]
En phase gazeuse, on peut travailler avec les \textbf{pressions partielles} ($K_p$).
\end{formuleb}

\begin{attentionb}[ce qui n'apparaît PAS dans $K$]
Les \textbf{solides purs} et \textbf{liquides purs} (dont le solvant) n'entrent \emph{pas} dans $K$. Seuls
les \textbf{gaz} et les \textbf{espèces dissoutes} $(aq)$ y figurent. Ex. : pour
$\ce{FeO}(s) + \ce{CO}(g) \rightleftharpoons \ce{Fe}(s) + \ce{CO2}(g)$, $\;K = \dfrac{[\ce{CO2}]}{[\ce{CO}]}$.
\end{attentionb}

\begin{theoremeb}[lire $K$ et de quoi il dépend]
$K \gg 1$ : équilibre déplacé vers les \textbf{produits} ; $K \ll 1$ : vers les \textbf{réactifs} ;
$K \approx 1$ : coexistence. \;
$K$ dépend \textbf{uniquement de $T$} ; il ne dépend \emph{pas} des concentrations initiales ni d'un
catalyseur. Multiplier les coefficients par $n$ élève $K$ à la puissance $n$ ; inverser l'équation donne $1/K$.
\end{theoremeb}

\section{Déplacement d'équilibre : loi de Le Chatelier}

\begin{theoremeb}[loi de Le Chatelier (modération)]
Si l'on perturbe un système à l'équilibre, il évolue dans le sens qui \textbf{s'oppose} à la perturbation.
\end{theoremeb}

\renewcommand{\arraystretch}{1.25}
\begin{center}\small
\begin{tabular}{@{}>{\raggedright\arraybackslash}m{4.6cm}>{\raggedright\arraybackslash}m{11.3cm}@{}}
\toprule
\textbf{Perturbation} & \textbf{Effet} \\
\midrule
\rowcolor{primary!6} $T \nearrow$ & déplacement dans le sens \textbf{endothermique} ; \emph{$K$ change} (endo : $K\nearrow$ ; exo : $K\searrow$). \\
$p \nearrow$ (gaz) & vers le côté qui a \textbf{moins} de moles de gaz (sans effet si autant de gaz des deux côtés). \\
\rowcolor{primary!6} ajout d'un \textbf{réactif} & vers la \textbf{droite} (le système le consomme) ; ajout d'un produit $\to$ vers la gauche. \\
ajout d'un \textbf{solide/liquide pur} & \textbf{aucun effet} (n'apparaît pas dans $K$). \\
\rowcolor{primary!6} \textbf{catalyseur} & \textbf{aucun déplacement} : accélère les deux sens, atteint l'équilibre plus vite (n'affecte pas $K$). \\
\bottomrule
\end{tabular}
\end{center}
\renewcommand{\arraystretch}{1.0}

\begin{attentionb}
La \textbf{pression} n'agit que sur les \textbf{gaz}, et seulement si leur nombre de moles \emph{change}
entre les deux côtés. Ajouter un gaz \textbf{inerte à volume constant} ne change pas les pressions
partielles : aucun effet. Seule la \textbf{température} modifie la valeur de $K$.
\end{attentionb}

\end{document}
